\chapter{Introduzione}
\paragraph{Introduzione al progetto e analisi delle specifiche} Il progetto si prefigge l'obiettivo di costruire una base di dati per una scuola di musica gestendo molteplici scenari: 
\begin{itemize}
    \item Gestione dei corsi individuali e collettivi 
    \item Gestione degli allievi
    \begin{itemize}
        \item Strumento suonato 
        \item Livello di preparazione  
        \item Docente assegnato
        \item Orario della lezione
    \end{itemize}
    \item Tracciamento delle presenze degli allievi 
    \item Esami di livello e saggi
\end{itemize}
Il database si dovrà inoltre occupare della gestione lato docente con: 
\begin{itemize}
    \item Specializzazioni docenti
    \item Disponibilità 
    \item Compensi
\end{itemize}
\newpage
Infine la parte commerciale e amministrativa con: 
\begin{itemize}
    \item Noleggio strumenti musicali 
    \item Vendita spariti ed accessori 
    \item Aule con strumentazione 
    \item Concerti degli allievi 
    \item Workshop con artisti esterni
\end{itemize}
\paragraph{Studenti e certificazioni}
Ogni studente può avere assegnato uno o più strumenti al momento dell'iscrizione nella scuola; di conseguenza, l'assegnazione a un docente per quello stesso strumento in base alla compatibilità degli orari. 
\\
Ogni lezione rappresenta uno slot di disponibilità per il docente, il tentativo di inserimento di una lezione all'interno di un slot già precedentemente occupato non è una operazione permessa all'interno del database. 
\\ 
Tutti gli studenti sono classificati in base al loro livello calcolato sulla base di esami conseguiti
\\
In base a ciò sarà possibile determinare quali esami possono effettuare gli studenti (e.g uno studente con livello 2 non può effettuare un esame di livello 4 senza prima aver ottenuto la certificazione per il livello 3) 
\paragraph{Tracciamento pagamenti}
La regolarità amministrativa viene tracciata attraverso il controllo dei regolari pagamenti delle lezioni attraverso una \texttt{VIEW} per visualizzare i pagamenti in ritardo
\newpage
\paragraph{Aule con strumentazione}
Le aule destinate alle lezioni presentano dotazioni strumentali eterogenee. In fase di prenotazione da parte di un docente, il sistema dovrà gestire due aspetti fondamentali:
\begin{itemize}
    \item \textbf{Compatibilità strumentazione}
    \begin{itemize}
        \item È necessario che l'aula scelta possieda l'equipaggiamento necessario per lo svolgimento della lezione
    \end{itemize}
    \item \textbf{Esclusività della prenotazione} 
    \begin{itemize}
        \item È necessario un vincolo di integrità per impedire sovrapposizioni durante la prenotazione, evitando che la medesima aula venga prenotata contemporaneamente da più docenti.
    \end{itemize}
\end{itemize}
\paragraph{Noleggio strumenti}
Il noleggio strumenti richiede l'introduzione di rigorosi vincoli temporali e di disponibilità al fine di prevenire l'incongruenza dei dati quali: 
\begin{itemize}
    \item \textbf{Coerenza temporale} Un vincolo di dominio (implementabile attraverso la clausola \texttt{CHECK}) per far si che la data di fine noleggio non possa essere antecedente a quella di inizio
    \item \textbf{Verifica della disponibilità} Un vincolo che impedisce il noleggio di uno strumento (identificato attraverso una chiave univoca) attualmente già prenotato (attraverso l'opportuno uso di \texttt{TRIGGER} al momento della prenotazione)
\end{itemize}
\paragraph{Vendita spariti e accessori}
Per quanto concerne il modello relativo alla vendita di articoli e accessori (corde, plettri, cavi, etc.) i vincoli si basano sull'integrità della gestione del magazzino e della correttezza delle transazioni, in particolare: 
\begin{itemize}
    \item \textbf{Validità dei domini} al momento dell'acquisto di un oggetto è necessario verificare che l'oggetto acquistato sia in una quantità $> 0$ e l'importo della transazione un numero positivo
    \item \textbf{Verifica delle giacenze} Le giacenze in magazzino devono essere gestite attraverso due vincoli fondamentali: 
    \begin{itemize}
        \item Al momento della vendita la quantità dell'oggetto acquistato deve essere pari o superiore a quella in giacenza all'interno del magazzino
        \item Una volta evaso l'ordine è necessario un \texttt{TRIGGER} che si occupi di decrementare gli oggetti in giacenza in base alla quantità ordinata 
    \end{itemize}
\end{itemize}
